\documentclass[a4paper,11pt, oneside]{book}
\usepackage[utf8]{inputenc}
\usepackage[francais]{babel}
\usepackage[T1]{fontenc}
\usepackage{graphicx}
\usepackage{float}
\usepackage{wrapfig}
\usepackage{setspace}
\usepackage{geometry}
\usepackage{hyperref}
\usepackage{multicol}
\usepackage{etoolbox}
\usepackage{color}
\usepackage[explicit,pagestyles]{titlesec}
\usepackage[absolute,overlay]{textpos}
\usepackage{fancyhdr}

\usepackage{eurosym}
\usepackage{titlesec}
\usepackage[]{algorithm2e}


% ====== CONFIG ========

%\setmainfont{Roboto Light}

%\setmainfont{Roboto Light}
%\setsansfont{Roboto}
%\setmonofont{Roboto}
%\newfontfamily\light{Roboto Slab Light}
\graphicspath{{img/}}
\setlength{\unitlength}{1mm}
\makeatletter

\definecolor{primary}{RGB}{44, 62, 80}


\titleformat{\chapter}[display]{\huge}{\thechapter \quad #1}{0pt}{}
\titlespacing{\chapter}{0pt}{0pt}{0pt}

\titleformat{\section}[display]{\LARGE}{}{0pt}{\thesection \quad #1}


\setlength{\TPHorizModule}{1mm}
\setlength{\TPVertModule}{1mm}
\def\sizeMedia{38}
\def\size{3.8cm}
\def\sizeMargin{0.2cm}
\def\margin{2}
\def\fixMargin{0}

\pagestyle{plain}

\author{Romain Cahu, Victor Cholley-Barroyer, Yann Prono, Quentin Tardivon}
\date{\today}

\def\school{TELECOM Nancy}
\def\schoolAddress{193 Avenue Paul Muller}
\def\schoolPostalCode{54602}
\def\schoolCity{Villers-lès-Nancy}
\def\schoolCodeAndCity{\schoolPostalCode, \schoolCity}
\def\schoolYear{2016 - 2017}

\def\appName{MALG Project}
\def\widthImage{1}
\def\todo{{\color{red}\Huge{TODO}}}

\def\schoolYear{2016 - 2017}

% ====== END CONFIG ========


\begin{document}

	\begin{titlepage}
		\thispagestyle{empty}

{\color{primary}


\includegraphics[width=4.0cm]{img/school-logo.eps}
\hspace{9mm}
\includegraphics[width=4.0cm]{img/collegium-logo.eps}
\hspace{5mm}
\includegraphics[width=4.0cm]{img/university-logo.eps}

\vspace{0.5cm}

	\begin{center}


			{\color[rgb]{0.8,0.8,.8}\rule{\textwidth}{0.8pt}}
			\vspace{0.2cm}

			\baselineskip=3pt
			{\huge \bfseries{Modèles et algorithmes}}\\
			\vspace{0.5cm}
			{\Large \bfseries{Tri à bulles - Algorithme P1 - Triangle de Floyd}}
			\vspace{0.5cm}

		{\color[rgb]{0.8,0.8,.8}\rule{\textwidth}{0.8pt}}

		\vspace{1cm}

		\Large{Groupe 10}\\

		\vspace{1cm}
		\includegraphics[width=0.4\textwidth]{logo.png}


		\vspace{2cm}

		\Large{Romain Cahu}\\
		\Large{Victor Cholley-Barroyer}\\
		\Large{Quentin Tardivon}\\
		\Large{Yann Prono}\\
		\vspace{1cm}
		\large{\schoolYear}
	\end{center}

}

	\end{titlepage}

	\newpage\newpage\null\thispagestyle{empty}
	\newpage
		\tableofcontents
		\thispagestyle{empty}

	\chapter{Introduction}
	\setcounter{page}{1}
		Lorem ipsum dolor sit amet, consectetur adipisicing elit, sed do eiusmod tempor incididunt ut labore et dolore magna aliqua. Ut enim ad minim veniam, quis nostrud exercitation ullamco laboris nisi ut aliquip ex ea commodo consequat. Duis aute irure dolor in reprehenderit in voluptate velit esse cillum dolore eu fugiat nulla pariatur. Excepteur sint occaecat cupidatat non proident, sunt in culpa qui officia deserunt mollit anim id est laborum.

	\chapter{Triangle de Floyd}

		\section{Description de l'algorithme}

		Le Triangle de Floyd prend en paramètre un entier n et imprime à l'écran
		n lignes contenant chacune 1 puis 2 puis 3 jusque n entier.


		\section{Implémentation et preuves PlusCal - TLA+}

			L'algorithme PlusCal correspondant est relativement simple. On a étiqueté
			un maximum d'étape de l'algorithme pour pouvoir y référer plus facilement
			et pour prouver un maximum de préconditions et d'invariants.

			On a d'abord vérifier les invariants de boucles. En effet, le fait d'être
			dans la boucle p2 (c'est à dire que le programme est au point p5) implique
			que i est plus petit que n. De la même façon, en p4 c est plus petit ou égal
			à n.

			On s'est ensuite assurer qu'il n'y avait pas d'overflow ou d'underflow au sein
			du programme. On a pour cela défini le max à 65535 correspondant à l'entier
			maximum réprésentable dans un système 16bits non signé.

			Enfin, la correction partielle du programme est assurée par la propriété
			Correctness. Lorsqu'on arrive en pc="Done", c'est à dire à la fin du programme
			on s'assure que le dernier entier imprimé correspond à la formule suivante:
			$a = [n \times (n + 1) / 2] + 1$

			Cette formule correspond à la somme suivante: $$\sum_{k=1}^{n}k$$

		\section{Implémentation et preuves Frama-C}

	\chapter{Tri à bulles}

		\section{Description de l'algorithme}

		Le tri bulle consiste à comparer répétitivement les éléments consécutifs d'un tableau,
		et à les permuter lorsqu'ils sont mal triés. Ici, il prend en paramètre un
		tableau d'entier à trier et le nombre d'élément de ce tableau.


		\section{Implémentation et preuves PlusCal - TLA+}

		De la même façon que le Triangle de Floyd l'implémentation en PlusCal est
		relativement rapide. La seule difficultée est la declaration de tableau.

		Ici on a fait la preuve de l'algorithme sur un tableau fixe de 4 éléments.

		De la même façon que le programme précédent on a d'abord vérifier les invariants
		de boucles dans les invariants Prec1 et Prec2.

		On a ensuite vérifier qu'après une opération de swap, les deux éléments
		du tableau étaient localement ordonnés dans la propriété Prec3. On également
		vérifier l'absence d'overflow et d'underflow de la même façon que pour le programme
		précédent.

		Enfin, on a vérifié la correction de l'algorithme en s'assurant que le tableau
		résultat était globalement trié. Pour cela on a déclaré une fonction isSorted vérifiant
		que pour tout élément du tableau on tableau[k] <= tableau[k+1].
		On a ensuite appliqué cette fonction à la fin de l'algorithme, c'est à dire
		lorsque pc = "Done".


		\section{Implémentation et preuves Frama-C}

	\chapter{Algorithme P1}

		\section{Description de l'algorithme}

		L'algorihme P1 était un algorithme inconnu. On s'est rendu compte lors de la
		preuve de ce dernier qu'il calcule les termes de la suite recursive de 2nd
		ordre suivante: $b_n = 2 \times b_{n-1} + 2 \times b_{n-2}$ avec $b_0 = 1$ et $b_1 = 1$.

		\section{Implémentation et preuves PlusCal - TLA+}

		De la même façon que les deux algorithmes précédents, l'implémentation PlusCal
		était relativement simple.

		Toujours comme les algorithmes précédents, on a vérifié les conditions de la boucle
		p9 ainsi que les conditions des tests p5 et p7 respectivement dans les propriétés
		Prec5, Prec1 et Prec3. On a également vérifier le non respect de ces conditions
		dans les blocs else de ces différentes conditions dans Prec2 et Prec4. On a
		également la correction partielle pour v = 0 et v = 1 dans la propriété Res,
		ce qui correspond aux termes $b_0$ et $b_1$ de la suite. On pouvait vérifier
		que les éléments de la suite était bien des entiers naturels dans la propriété
		PartialCorrectness. On a également, comme pour les algoritmes précédents vérifier
		l'absence d'overflow et d'underflow.

		En revanche, la correction totale de l'algorithme pose problème. En effet,
		on a effectué la résolution de l'équation caractéristique de la suite pour trouver
		l'expression de son terme général en fonction de $b_0$ et $b_1$.
		On trouve après simplification:
		$$1/2 \times [(1 + \sqrt 3)^n + (1 - \sqrt 3)^n]$$

		Or, TLA+ ne nous permet de réaliser de racine carré ou de puissance avec
		un réel ($1/2$ par exemple).

		On peut en revanche imaginer une version de l'algorithme qui calculerait
		directement le terme n de la suite sans utilisation de boucle for, ce qui
		augmenterait la vitesse de l'algorithme pour des grandes valeurs de n.


		\section{Implémentation et preuves Frama-C}

	\chapter{Conclusion}
		Lorem ipsum dolor sit amet, consectetur adipisicing elit, sed do eiusmod tempor incididunt ut labore et dolore magna aliqua. Ut enim ad minim veniam, quis nostrud exercitation ullamco laboris nisi ut aliquip ex ea commodo consequat. Duis aute irure dolor in reprehenderit in voluptate velit esse cillum dolore eu fugiat nulla pariatur. Excepteur sint occaecat cupidatat non proident, sunt in culpa qui officia deserunt mollit anim id est laborum.

\end{document}
